\section{Discussion}\label{sec:discussion}
The experiments connect the requested event, checked program result, and
credited execution. They expose admission distortion, recovery requirements,
and unsupported attribution under separate contracts. Their contribution is
complementary empirical evidence rather than repeated estimates of one effect.

\paragraph{RQ1: a measured distortion with a tested mechanism}
Matched instances fix the recovery problem; regrading fixes each program and
oracle outcome. The earlier filter is narrower than the H1-R3 prompt, so its
$-32.0$ point movement measures a contract-mismatched rule on this cohort.
Representation-associated \texttt{hashlib} imports explain most of it, and
output-preserving cleanup removes the gap. This identifies the score-changing
decision and its source trigger, not disagreement between compliant policies.

\paragraph{RQ2: feasibility, construction, and consequence}
The fixed executor establishes operation feasibility; V4 tests construction.
Parameter errors and reduced-information variants locate recovery
requirements. AES-CTR demonstrates supplied-chain payload release, while
C2b and wire editing distinguish policy enforcement from an unverified
handshake. Each target predicate defines the supported consequence.

\paragraph{RQ3: authenticity and consistency are complementary}
E5's committed candidate match leaves acquisition unresolved. The attribution
studies test explicit evidence requirements rather than reconstruct that
historical acquisition path.
The HEB-1 and IAR-B/IAR-C cases distinguish candidate correctness from
execution attribution. Join-only and signature-only checks each promote
unsupported claims that combined checks reject. The join validator's
rejection of honest partial claims reflects its full-attribution-only
contract. The supplement retains each constructed corpus's denominator.

A serialization crash and interface-contradicting source policy each affected
72/72 records, invalidating that campaign. Together with H1 V2's scorer
mismatch, these diagnostics show why failure-stage retention is necessary
to interpret a terminal outcome.

\subsection{Contribution to Security-Agent Measurement}
ABC and AgentSuite analyze benchmark validity across task, environment, and
evaluation components \cite{abc,agentsuite}; cryptographic-capability work
supplies the task context \cite{aicrypto,cryptanalysisbench}.
Our empirical contribution combines program-level contrasts, recovery
diagnostics, and tested execution-claim policies. H1-R3 quantifies admission
distortion; target studies check consequences; attribution cases test when
correct or authentic evidence supports only a weaker claim. The claims
connect the experiments while their populations remain separate.

\subsection{Coverage Across Settings}\label{sec:external-coverage}
The completed studies broaden different claims along different axes
(Table~\ref{tab:external-coverage}). V4's capture-first loop and the
post-freeze OpenCode CLI studies retain accepted constructions under two
interfaces and model labels. Their differing tasks, observations, and budgets
support feasibility across settings, not a matched harness or model effect.

\begin{table}[t]
\centering\footnotesize
\caption{Existing evidence and the scope it supports. Rows retain separate
populations; diversity of settings does not replicate the admission effect.}
\label{tab:external-coverage}
\begin{tabularx}{\columnwidth}{@{}P{1.45cm}Y@{}}
\toprule
\tablehead Claim & Coverage and boundary\\
\midrule
Admission & Two H1-R3 families; one route and policy pair. The 32-point
change has no independent-benchmark replication.\\
Construction & V4 and OpenCode CLI establish recovery or supplied-chain
execution under distinct local contracts, without a pooled rate.\\
Held-out & Two development and two held-out cells share two secrets; the
alphabet shift tests within-family transfer.\\
Attribution & IAR-B/IAR-C vary case construction; configured in-toto tests
another implementation of the same policy.\\
Public releases & Six external releases supply documentary reconstruction
evidence, not executable validation.\\
\bottomrule
\end{tabularx}
\end{table}

IAR-C's separate generator and adjudicator share the schema and cryptographic
module, and its preserved cases were used in calibration. The in-toto
inspection implements the same policy. These extend constructor and
implementation coverage; the admission magnitude remains local.

\subsection{Implications for Agent Evaluation}
\paragraph{Publish the contract that the executor enforces}
Publish imports, runtime, observations, feedback, budget, and predicate.
Check prompt--executor consistency; distinguish fixed-program regrading from
new generation when a rule changes.

\paragraph{Retain the decision path}
Keep source, runtime, predicate, provider, and acceptance states distinct.
Handshakes, operation feasibility, and agent completion require their own
endpoints and denominators.

\paragraph{Bind acceptance to the claimed execution}
Retain candidate-to-verdict bindings, trusted observer/reference identities,
and replay policy. Attribution requires authenticated, claim-consistent
evidence; acquisition, host isolation, and immutable hosted-model identity
remain separate requirements.

The supplement maps each claim to its retained records and separates
recomputable results from documentary or recorded-only evidence.
